\section{Conclusion}\label{sec:conclusion}

The shortfall of the Federal Reserve's mortgage runoff beneath its Quantitative Tightening caps was not solely a product of high interest rates. It followed from the structural design of the U.S. mortgage contract interacting with a prepayment environment mostly rate-inelastic relative to the caps. Under the production floor form the $\beta_1 = 0$ null recovers 85.7\% of the benchmark at the headline calibration (35.6\% under the additive form). So the caps could not have bound under any rate response above the baseline turnover the null embeds. How much of the remainder sits outside household-level decision-making is a question the cross-design test of Section~\ref{sec:robustness-crossdesign} leaves open. This paper's central empirical finding is the rate-inelastic decomposition, established in Section~\ref{sec:hazard} on real Freddie Mac loan-level data. The agent-based model of Section~\ref{sec:abm} runs on the synthetic household population specified in Section~\ref{sec:method}. It stands beside that finding as a stress-test companion, since Section~\ref{sec:robustness-crossdesign} confines what its verdict can mean. The two are complementary, and neither is derived from the other. They are not independent in calibration, though: both are anchored to a baseline turnover floor read off discount-cohort turnover, and at the in-sample calibration point to the same 2023--24 read (Section~\ref{sec:pathb}).

Rational, behaviorally extended, or subjected to a direct burnout test, the ABM cannot explain the majority of the Federal Reserve's cap shortfall. A structurally distinct hazard framework is built entirely from loan-level survival dynamics, with no household utility function at all. In its microsimulation form, on the benchmark-consistent accounting basis, it recovers 91.3\% of the full empirical benchmark at the off-window floor this paper headlines. It recovers 88.2--93.7\% across that floor's defensible range and 97.9\% at the in-sample calibration point. So the miss is 8.7 points on the headline calibration and 2.1 points on the in-sample one. That is a level-accuracy distinction, not a timing one. No estimator's monthly co-movement survives detrending (Section~\ref{sec:pathb}). The estimated cohort-hazard form lands at a comparable aggregate level but does not corroborate it. Its own coefficient's sign is not statistically distinguishable from zero. It carries no valid production-spec uncertainty and is excluded from headline figures (Section~\ref{sec:patha}).

The composition of that recovery disciplines what the contrast can mean. Under the production floor form the $\beta_1 = 0$ null recovers 85.7\% on the same basis at the headline floor (35.6\% under the additive form) and 88.7\% in-sample. The paper's own ablations bound the named loan-level mechanisms: burnout $0.45$ points at the headline floor and $0.86$ in-sample, FICO/LTV covariates $1.3$ points, the delinquency pipeline \$0.39 billion. No servicer channel beyond that pipeline is modeled at all. What delivers the level is scheduled amortization, the seasoning baseline, and the turnover floor. The shortfall against the phased caps is mostly rate-inelastic, because the caps sat far above what any plausible prepayment environment would have delivered. But against the Federal Reserve's own ex-ante projection rather than the non-binding cap, the lock-in channel accounts for roughly half the genuine surprise under the central allocation. That share runs 23--49\% at the off-window floor and 38--80\% at the in-sample calibration point across the two disclosed intra-2022 allocations (Section~\ref{sec:method-benchmark}). The elasticity's contribution is a bounded margin over that rate-inelastic base: $+2.3$ to $+9.1$ points under the production floor form ($+9.2$ at the in-sample calibration point; $+11.2$ under the competing-risks form). Inside that interval, $+5.6$ is the value at the mid-grid anchor. Measured corrections to the floor and the accounting basis move that value in both directions within the interval (Section~\ref{sec:identification}). The cross-design answer is asymmetric. The hazard framework's aggregate recovery survives a fully synthetic population. The ABM's share moves with calibration choice and data source: 59.3\% recalibrated, 76.3\% at the book's composition. That range sits above the pre-committed threshold. Above it, Section~\ref{sec:robustness-crossdesign} concludes the contrast cannot, on that test alone, be cleanly attributed to modeling paradigm. The path-level timing question remains open for every estimator.

Section~\ref{sec:abm-danish}'s Danish counterfactual retires a superseded figure. Recalibrated against \citepos{berger2026} separately-estimated Danish elasticities (unrefereed working paper, draft July 2026), the \$925.5 billion mechanism-extrapolated gap collapses by an order of magnitude. Anchoring the transplant so that it changes only the payoff rule---the anchor consistent with the baseline turnover floor's semantics---signs it at $+\$61.2$ billion at the band's zero-refinance edge, 8\% of the benchmark and the image of the lock-in marginal on the in-sample calibration both legs share. At the off-window floor the paper headlines, the rule-only gap is $+\$28.2$ billion, 34\% below that marginal (Section~\ref{sec:abm-danish}). The sign is anchored, not estimated: pinning the Danish leg at the U.S. zero-gap hazard guarantees it against a book that is out of the money almost everywhere. The bracketing anchor that also imports Danish baseline mobility reverses it to $-\$99.9$ billion, on a Danish leg whose 3.39\% mean prepayment sits below the turnover floor the U.S. legs impose; that inconsistency demoted it to a bracket. The payoff-rule question therefore ends at a modest magnitude and a stated anchor. One qualification is priced as a bracket of its own: the Danish leg credits retired balance at par, and the pre-committed bracket (runs \texttt{buyback\_credit\_bracket} and \texttt{buyback\_discount\_rederived}) prices both incidence readings---the gap stands at $+\$61.2$ billion under the balance-adjustment reading and reverses to $-\$51.0$ billion under the cash-haircut reading at the re-derived discount. The relief therefore stands on the benchmark's own face denomination, while as reserve drain the rule costs the Federal Reserve exactly the par windfall it would otherwise have extracted from moving households. On that reading, legs (2) and (3) of the trade-off are the same dollars seen from two sides. The mobility relief such a rule would provide remains the better-supported and larger effect.

The remaining leg of the trade-off is secondary-market liquidity. The Danish system relies on strict one-to-one match funding (the Balance Principle); the United States relies on the To-Be-Announced (TBA) forward market, which requires highly homogenized, pooled mortgages. \citet{vickery2013} show that this homogeneity produces a measurable liquidity premium: TBA-eligible pools trade at significantly lower cost than comparable ineligible specified pools. \citet{gao2017} measure the same differential directly in transaction data, and \citet{song2019} document the dollar-roll market that finances TBA positions. \citet{campbell2013} frames the U.S.--Danish contrast as a mortgage-design problem in exactly these terms. A market-value repurchase feature prices each payoff at the market value of the funding bond series (a series-level price common to every borrower funded by that series on a given day). What it strains is therefore not homogeneity itself but its \emph{locus}: Denmark standardizes at the level of large, tap-issued bond series, the TBA market at the level of the deliverable pool. The Danish callable market's own depth shows that market-value redemption and liquid secondary trading coexist under the series convention \citep{berg2018}. The binding cost for the United States is migrating its issuance architecture from pool-level to series-level standardization, not a liquidity impossibility. Adopting the rule without that redesign would nonetheless risk a documented source of liquidity in the transition. A hedging-side literature sharpens what the two payoff rules allocate. Prepayment risk is borne and priced by specialized MBS investors, not by diversified households \citep{gabaix2007}. Its aggregate hedging feeds back into the level and volatility of long-term rates \citep{malkhozov2016}. The behavioral component is the part that cannot be hedged away, only bounded \citep{perotti2024}. The Danish rule does not remove that behavioral residual --- the borrower retains the par call on the rate-decline side --- but it removes the rate-rise extension leg on which this cycle's exposure accrued. The choice between payoff rules is thus also a choice about which side of the convexity profile the imperfectly hedgeable risk sits on.

U.S. policymakers weigh three considerations: (1) deep, homogenized secondary market liquidity (the TBA market); (2) continuous household mobility (avoiding the lock-in effect); and (3) convexity-neutral duration profiles (predictable institutional cash flows). This paper's evidence is that they do not form the strict trilemma a three-way ``trade-off'' would imply. That reading holds under its production floor form; the additive form roughly doubles the margin at issue. The evidence of Sections~\ref{sec:abm} and~\ref{sec:hazard} is that the U.S. has prioritized (1) at real cost to (2): household mobility is genuinely constrained, while the cash-flow shortfall against the caps is mostly rate-inelastic. The Federal Reserve staff's own record and the ex-ante projection corroborate that composition as composition, not as numerical agreement, and not independently of it, since the projection may already have embedded part of the channel (Section~\ref{sec:lit}). Within that composition, the lock-in channel contributes an identified margin of $+2.3$ to $+9.1$ points at the off-window floor, $+5.6$ at the mid-grid anchor inside it ($+9.2$ at the in-sample calibration point). That constraint is not evenly borne. The marginal's per-balance intensity tilts toward below-740-FICO ($1.23\times$--$1.26\times$) and above-80-LTV ($1.23\times$) borrowers (Section~\ref{sec:identification}), and the statutory assumability carve-out that would relieve it is bounded above by the 20.4\% Ginnie share, with take-up unmeasured; the mobility cost therefore falls disproportionately on the households with the fewest refinancing alternatives. What I no longer find support for is the assumption that (2) and (3) trade off sharply under face accounting, in which fixing mobility would cost face-value cash flow. Under the rule-only transplant, Section~\ref{sec:abm-danish}'s recalibrated Danish counterfactual finds that a market-value repurchase rule would have relieved (2), while its effect on (3) is incidence-conditional: an $\approx\$61$ billion improvement under face accounting, a $-\$51.0$ billion cost under cash accounting at the re-derived discount (the foregone household-to-bondholder transfer). That reading inherits the marginal's floor-and-band conditionality and its statement at the in-sample calibration point (Section~\ref{sec:identification}). The dissolution rests on thinner ground than the framing suggests, so I state it plainly. The rule-only Danish gap is the lock-in marginal viewed from the counterfactual side, so ``(2) and (3) do not trade off sharply'' reduces to ``the marginal is small''. Under the production floor form the marginal is bounded small by construction. By contrast, under the disclosed additive form, which never censors the elasticity, the same counterfactual is about $+11$ points and nearly floor-invariant (Section~\ref{sec:robustness-floor}). Because \eqref{eq:pathB} combines the turnover floor with the voluntary hazard by a hard maximum, switching the elasticity off can raise the hazard only by the excess of the $\beta_1 = 0$ hazard over the floor---and not at all where the floor binds on that leg itself, 14.3\% of loan-months at the in-sample calibration point (Section~\ref{sec:identification}). The dissolution is therefore a statement about a calibrated ceiling, not an independent finding, and it is form-conditional. What survives across forms is that the trade-off is modest, not that it is absent. On either form, the binding constraint against adopting such a rule is not a mobility--cash-flow trade-off but the TBA liquidity premium described above. That is an argument from the cited literature, not a result of this design: nothing here prices a series-level call, so the investor-side cost of the transplant is named and left unpriced (\citet{diep2021} show that prepayment risk is compensated in the cross section of MBS returns, which is the price this design does not measure). This is the two-sided exchange the paper opened with: (1) purchased in the currency of (2), with a small side payment in (3) whose sign depends on the buyback's incidence. As long as the U.S. maintains the par-payoff MBS structure within the TBA framework, severe duration extension should be expected to recur in future tightening cycles that begin, as this one did, with a large outstanding stock of deeply in-the-money, below-market-coupon mortgages. The state dependence documented by \citet{berger2021} and \citet{eichenbaum2022} makes the extension conditional on that stock, not on tightening per se. The recurrence claim is itself an extrapolation beyond what this design tests. The accounting machinery is also episode-shaped: with essentially the whole book out of the money, the marginal's sign is forced and the refinancing margin is inert, so a mixed-moneyness window would require re-deriving both. This design tests one episode. The duration extension itself, however, is mostly the arithmetic of scheduled amortization and baseline turnover against caps set above any plausible prepayment path; an alternative payoff rule would have trimmed it by the lock-in margin, not eliminated it, and would carry the liquidity cost above.

\subsection{Limitations and Avenues for Future Research}\label{sec:limitations}

Several limitations remain on both sides of the framework.

\textbf{Housing supply as a competing channel}: The mobility margin this paper imports is estimated in an environment of historically tight for-sale inventory. Housing supply enters here only as a friction scaler on the moving decision, never as a channel in its own right. A household that does not move because nothing is listed is observationally identical, in these data, to one that does not move because its coupon is below market. This design does not separate them. The direction is argued rather than measured. A binding supply constraint would absorb some of what is attributed to lock-in. That is why the scaled-null exercise that removes the non-rate share of the moving decision is reported alongside the headline rather than as a robustness footnote (Section~\ref{sec:identification}; \citealp{aladangady2024,hsieh2019}). No inventory series is used to price it.

\textbf{Distributional silence}: This paper measures a book-wide institutional cost and never asks who bears it. Ginnie Mae borrowers---20.4\% of SOMA face, and the population whose structurally faster speeds Section~\ref{sec:pathb} documents---sit outside the Freddie Mac sample entirely. So nothing here establishes whether the mobility cost falls differently on FHA and VA households than on the conventional borrowers the sample contains. I record that as a distributional limitation of the design, not merely as a coverage bound on the estimate.

\textbf{Sample scale and comparability}: The ABM's 10,000-household synthetic population and the hazard framework's 75,000-loan real Freddie Mac sample differ in both size and data source. Section~\ref{sec:robustness-scale} settles the size question on the production specification: population size changes precision, not central tendency. The completed cross-design pair of Section~\ref{sec:robustness-crossdesign} addresses the data-source question asymmetrically. The hazard framework's aggregate recovery is largely data-source-independent --- the ABM's is not.

\textbf{Data:} A dedicated ingest pipeline integrates the Freddie Mac 2017--2021 loan-level data (20 origination quarters) that both hazard paths use. Where real data files are unavailable, results are validated against a synthetic fixture before real-data estimation. Fifteen-year MBS, approximately 9\% of SOMA face value, remain excluded from cohort modeling on the hazard side only; the ABM side folds them in (Appendix~\ref{sec:robustness-pipeline}).

\textbf{Model:} On the ABM side, the monthly CPR path fit remains weak (raw $R^2$ of $-6.984$), despite a directionally sensible level response to rates. The production Danish counterfactual is a rule-only transplant onto U.S. households (Section~\ref{sec:abm-danish}): it does not estimate how Danish institutional frictions would themselves migrate. The Danish-level bracketing anchor shows the gap's sign is not robust to bundling baseline-mobility differences into the transplant. The debt-to-income check remains front-end only, without a measure of total household debt. On the hazard side, Path A's Poisson pseudo-maximum-likelihood estimator requires a mild ridge penalty to remain well-conditioned at 295 stratum fixed effects. That is a stabilization device: Section~\ref{sec:patha} shows the prior-spec (v3) coefficients equal the unpenalized fit on the estimable strata. Its holdout RMSE (spec v3: approximately 38 percentage points unweighted, 2.5 points exposure-weighted; the frozen spec v4 fit re-scored on the same holdout: 37.9 and 3.0), together with a lag-0 CPR correlation of $-0.378$, more negative even than the ABM's $-0.318$, confines its evidential contribution to the aggregate dollar level, as Section~\ref{sec:patha} details. Path B's literature-calibrated hazard magnitudes are directionally and quantitatively closer to the empirical benchmark than any ABM specification. But they are simulated on the 75,000-loan real Freddie Mac sample with elasticities calibrated from secondary literature rather than estimated jointly with the empirical benchmark itself. Path B's recovery percentage is reported across the calibration band (105.9\%--108.2\%), and now also with loan-level sampling uncertainty from an ex-ante-specified stratified bootstrap over the 75,000-loan draw (Section~\ref{sec:pathb}). An external second-agency replication (Fannie Mae loan-level data, same pipeline end-to-end) reproduces the lock-in marginal at $+8.68$ points against Freddie's $+9.20$, inside its pre-committed envelope. That is a sign-and-envelope replication, not a hazard-level match, and one whose envelope gate is 11.06 points wide. So the 0.52-point agreement, rather than the gate, is what should be read from it (Section~\ref{sec:pathb}). One global fact about the design, each of whose components is disclosed only locally: no outcome-holdout months exist anywhere in this paper. The nearest construction is the marginal's temporal holdout (Section~\ref{sec:robustness-floor}). There the turnover floor is calibrated on the window's first nineteen months, and the central-minus-null difference is then summed over the remaining twenty-three alone---months that enter that calibration nowhere. It returns $+\$45.1$ billion. That construction holds for the marginal, not for the level. And it is an input-stability check rather than an outcome validation: realized prepayment from the held-out months enters neither simulated leg, so it cannot falsify the marginal against data. The statement above stands without exception. The production baseline turnover floor is still calibrated on 2023--24 turnover inside the evaluation window. The pre-committed off-window re-measurement of Section~\ref{sec:robustness-floor} finds the defensible off-window level higher than 4\%, so the direction of any anchoring error shrinks rather than inflates the identified marginal; the headline marginal is now reported at that off-window level rather than the in-window one. Path A trains on 19 of the 42 QT months, and its ridge-selection holdout doubles as its out-of-sample diagnostic. Path B's \emph{level}, as distinct from its marginal, is still set by a floor measured in the window on which it is evaluated. Every recovery percentage in the paper is accordingly an in-sample or calibrated quantity in the months dimension. Ex-ante-specified sweeps, nulls, parity gates, and bootstraps discipline it, not clean holdout months. All headline figures are frozen in tagged run manifests, with reproduction details recorded in the replication package's run ledger. The marginal itself, at either floor, is a model-based decomposition conditional on the disclosed hazard structure (Section~\ref{sec:identification}). Causal identification in the econometric sense would require exogenous variation in rate gaps across otherwise-comparable loans. This design does not exploit it.

The Danish-regime recalibration (Section~\ref{sec:abm-danish}) calibrates both frameworks' Danish CPR response against \citepos{berger2026} separately-estimated moving and refinancing elasticities rather than an extrapolated U.S. mechanism, and it anchors the rule-only transplant at the U.S. zero-gap intercept. That collapses the institutional gap by an order of magnitude, to a signed $+\$61.2$ billion. That source is an unrefereed working paper, draft dated July 2026, and it is moving: its general-equilibrium refinance-in-place estimate rose roughly twentyfold between its January and July 2026 drafts (Section~\ref{sec:abm-danish}). Path B's full 5.5\%--7.7\% elasticity band is propagated through the microsimulation (Section~\ref{sec:pathb}). That exercise also surfaced and corrected a units error in the rate-gap coefficient's application (Appendix~\ref{sec:robustness-pipeline}). The symmetric companion cross-design test is complete (Section~\ref{sec:robustness-synthesis}). The hazard framework's survival structure recovers 106.0\% of the benchmark on a fully synthetic population, closely matching its real-data recovery, while its path-level timing does not transfer to synthetic data. Three questions remain open. The first is whether to add a seasonal multiplier to Path B's baseline hazard (Section~\ref{sec:pathb}). The second is a total-debt measure, or threshold sensitivity, for the ABM's front-end DTI gate (Appendix~\ref{sec:method-frictions}). The third is the monthly-frequency question the timing diagnosis leaves: the empirical path's variation at that frequency is seasonal, not rate-driven, and no estimator here models it. The second-agency replication is complete (Section~\ref{sec:pathb}; the envelope caveat is in the Model paragraph above). The cyclical-floor variant is likewise complete. What its grid measures is floor dispersion under the hard maximum rather than cyclicality, with the surviving co-movement component small and bounded (Section~\ref{sec:pathb}; Table~\ref{tab:uncertainty} notes; Appendix~\ref{sec:robustness-seasonalfloor}). The Ginnie composition bound has been upgraded from two published snapshots to a time-varying overlay built from the published monthly CPR series. That overlay leaves the conventional-share marginal at $+7.3$ points at the in-sample calibration point ($+4.4$ composed with the off-window headline floor), positive under every variant, with the Ginnie-specific level contribution inside the static bound (Section~\ref{sec:pathb}). Replication alone could not retire the vintage residual: the 2022 and pre-2017 vintages (23.1\% and 10.6\% of book face) lie outside the 2017--2021 universe of both agencies' samples used here. The observed-speed bound of Section~\ref{sec:pathb} now disciplines that extrapolation directly (\$11.7 billion, 1.5\% of benchmark, signed toward overstating trapped liquidity). What remains open on this dimension is therefore no longer the absence of a bound but the bound's own proxy assumption: observed Fannie segment speeds standing in for the SOMA book's, the same single-agency posture the production calibration itself carries. A modeled overlay in place of that proxy is not constructible under the committed sampler: the staged replication sample enforces the 2017--2021 vintage universe, the pre-QT snapshot rule leaves second-half-2022 originations with no admissible observation month, and the acquisition window closes at 2022Q4, so 2022 originations acquired in 2023 were never staged. The 2022 leg of the observed-speed bound (\$3.3 billion of the \$11.7 billion) therefore stands as that vintage's discipline. Appendix~\ref{app:composition} collects the formal composition comparison behind these bounds: sample versus book by agency, vintage, coupon, and term, with a cross-agency stability check on the dimensions the book's CUSIP-level disclosure cannot supply.

Sections~\ref{sec:abm} and~\ref{sec:hazard} establish that contrast under the synthetic household population this paper specifies. Section~\ref{sec:robustness}'s cross-design pair leaves the attribution asymmetric, and it inherits Section~\ref{sec:robustness-crossdesign}'s own hedge. The hazard framework's recovery survives a fully synthetic population (106.0\% versus 107.0\%). That is evidence its survival structure, not its data source, carries the aggregate level. But the recalibrated real-covariate ABM's 59.3\%, and its 60.2\% fifty-seed mean, cross the pre-committed threshold: above it the contrast cannot, on that test alone, be cleanly attributed to modeling paradigm rather than to calibration and data source. The path-level timing question, meanwhile, remains open for every estimator. That result cuts both ways. Recovering 106\% on a fully synthetic population shows that the calibration constants, not real loan-level structure, deliver the level, which is consistent with the rate-inelastic decomposition above. On the paradigm question itself I make no claim: the reading was pre-committed to a threshold, the threshold was crossed, and Section~\ref{sec:robustness-crossdesign} withdraws it. The decomposition this paper does deliver never rested on it.

\clearpage
\bibliographystyle{plainnat}
\bibliography{references}

\clearpage